\documentclass[tikz,border=6pt]{standalone}
\usepackage{ctex}
\usepackage{amsmath}
\usepackage{pgfplots}
\pgfplotsset{compat=1.18}
\usetikzlibrary{calc,angles,quotes,arrows.meta,positioning,decorations.markings}
\begin{document}
% 参考角 alpha 在四象限的对称：theta, pi-theta, pi+theta, -theta(=2pi-theta)
\begin{tikzpicture}[scale=4]
  \def\a{35} % 参考角
  \pgfmathsetmacro{\cx}{cos(\a)}
  \pgfmathsetmacro{\sx}{sin(\a)}

  % 坐标轴
  \draw[->,gray!70,thick] (-1.5,0) -- (1.6,0) node[right,black]{$x$};
  \draw[->,gray!70,thick] (0,-1.5) -- (0,1.6) node[above,black]{$y$};

  % 单位圆
  \draw[blue!60!black,very thick] (0,0) circle (1);

  % 四个对称点
  \coordinate (P1) at (\a:1);        % theta
  \coordinate (P2) at (180-\a:1);    % pi - theta
  \coordinate (P3) at (180+\a:1);    % pi + theta
  \coordinate (P4) at (-\a:1);       % -theta

  % 终边
  \foreach \P in {P1,P2,P3,P4}{ \draw[gray!55,thick] (0,0) -- (\P); }

  % 对称虚线：水平、竖直、对角
  \draw[teal!60,dashed] (P1) -- (P4);  % 关于 x 轴对称（竖直连线）
  \draw[teal!60,dashed] (P1) -- (P2);  % 关于 y 轴对称（水平连线）
  \draw[teal!60,dashed] (P1) -- (P3);  % 关于原点对称（过原点）

  % 点
  \foreach \P in {P1,P2,P3,P4}{ \fill[red!80!black] (\P) circle (0.018); }

  % 参考角弧（每个点处标出与 x 轴的参考角）
  \draw[purple!70,thick] (0.24,0) arc (0:\a:0.24);
  \node[purple!70,font=\scriptsize] at (\a/2:0.34) {$\alpha$};

  % 标签：角与坐标符号
  \node[anchor=south west,font=\scriptsize,align=left] at (P1)
    {$\theta=\alpha$\\ $(+\cos,\,+\sin)$\,I};
  \node[anchor=south east,font=\scriptsize,align=right] at (P2)
    {$\pi-\alpha$\\ II\,$(-\cos,\,+\sin)$};
  \node[anchor=north east,font=\scriptsize,align=right] at (P3)
    {$\pi+\alpha$\\ III\,$(-\cos,\,-\sin)$};
  \node[anchor=north west,font=\scriptsize,align=left] at (P4)
    {$-\alpha$（$2\pi-\alpha$）\\ $(+\cos,\,-\sin)$\,IV};

  % 说明框（半透明，放在左上空白）
  \node[draw=blue!50!black,fill=white,fill opacity=0.85,text opacity=1,
        rounded corners,font=\scriptsize,align=left,anchor=north west] at (-1.55,1.5)
    {以参考角 $\alpha$ 为基准的对称：\\[1pt]
     关于 $y$ 轴 $\to\ \pi-\alpha$\\
     关于原点 $\to\ \pi+\alpha$\\
     关于 $x$ 轴 $\to\ -\alpha$\\[2pt]
     $\sin,\cos$ 大小相同，\\ 符号由象限决定};
\end{tikzpicture}
\end{document}
